\MajorSection{الملاحظة الثانية}

\MinorSection{في القصيدة}

كلمات قليلة عن القصيدة نفسها. يبدأ حاجنا بإقامة المشهد، ويودّع القافلة المنطلقة إلى مكة. ويرى «ذيل الذئب»، دمِ گرگ، \textenglish{Dum-i-gurg}، أو ضوء الذئب، \textenglish{lykaugés}، أو أول الفجر، \textenglish{Diluculum}؛ وهو ضوء الفجر البروجي، أول ضربات البياض الخافتة المشعّة من أسفل الأفق الشرقي. ويصاحبه نفَس الصباح، دمِ صبح، \textenglish{Dam-i-Subh}: تيار هواء يكاد لا يُدرك إلا بازدياد البرد، ويفترض علماء وظائف الأعضاء المسلمون أنه الصلاة المبكرة التي ترفعها الطبيعة إلى العلة الأولى. ومن الواضح أن غولِ بيابان، شيطان الصحراء، تجسيد لمخاوف الإنسان وللأخطار المحيطة بالسفر في البراري. والبرية التي لا يسكنها إلا هو، الله، هي قفر عظيم مروّع، دشتِ لا سوى هو؛ وجبل الله المقدس هو عرفات قرب مكة، الذي تبلغه القافلة بعد مرورها بالمدينة. وينتهي القسم الأول بنواح موجع، إذ إن «لقاءات هذا العالم تحدث على طريق الفراق»؛ ويحتوي الأصل أيضًا:

\begin{SourceQuotation}
يخدّر برد الحزن فكري؛ وأحسبني أسمع ناقوس الرحيل؛\\
فيما يموت وراء ذلك الخط الأزرق الرقيق رنين جرس الجمل.
\end{SourceQuotation}

ويقتبس القسم التالي الوجوه المتنوعة التي ظهرت بها الحياة لمعلّمي البشرية الحكماء والحمقى. ويأتي حافظ أولًا؛ فتُقتبس أبياته المعروفة التي تبدأ بـ«شبِ تاریک و بیمِ موج»، أي الليل المظلم وخوف الموج، إلى آخرها. والحور جمع أحور، وبالكامل أحور العين: فتاة يشتد بياض عينيها حيث ينبغي أن تكونا بيضاوين، ويشتد السواد في سائر مواضعهما؛ ومن هنا كلمتنا السخيفة \textenglish{Houries}. ثم يأتي عمر الخيام، الذي أضفى معنى روحيًا على التصوف، كما أضفى الصوفية، أي العارفون، معنى روحيًا على التزمت الإسلامي. والأبيات المشار إليها هي:

\begin{SourceQuotation}
تعرفون، يا أصدقائي، بأي احتفال شجاع بالشراب أقمت في بيتي زواجًا ثانيًا؛\\
طلّقت العقل العجوز العقيم من فراشي، واتخذت بنت الكرمة زوجة.\\
المقطع ٦٠ من ترجمة السيد فيتزجيرالد.
\end{SourceQuotation}

تُستعمل «الخمر» هنا بمعناها الصوفي: الحب المستغرق لنفس النفوس. وكان عمر مكروهًا مخوفًا لأنه تكلم بجرأة، بينما لجأ إخوته الصوفية إلى التلميح. وقد صيغ اقتباس ثالث ليشبه «نشيد الحياة»، رغم الابتذال والعادية المؤلمة اللذين يميزان المدرسة الإنجليزية الأمريكية المتشبّهة بشيلر. وفُعل الشيء نفسه بكلمات عيسى، أي يسوع؛ فالمؤلف الواسع الاطلاع على الإنجيل أراد الإشارة إليه بوضوح. ورُجم منصور الحلاج، أي منظف القطن، لأنه نطق نطقًا فجًّا بعقيدة وحدة الوجود: «أنا الحق»، أي الله، و«ليس في جبتي إلا الله». ورسم دمه على الأرض العبارة الأولى. وأخيرًا يأتي اقتباس من «سردانابالوس ابن أناكينداراكس»، إلى آخره؛ وقد تعني الكلمة اليونانية \textenglish{paîze} اللهو، لكن السياق يحدد نوع اللهو المراد. والزاهد هو المؤمن حرفيًا بنص الشريعة، في مقابل الصوفي الذي يؤمن بروحها؛ ولذلك يُدعى الأول ظاهريًا، أي خارجيًا، والآخر باطنيًا، أي داخليًا. ويُقتبس موسى لأنه تجاهل الثواب والعقاب المقبلين. أما «الأبديتان»، فيقسم الميتافيزيقيون الفرس والعرب الأبدية، أي نفي الزمان، إلى نصفين: الأزل، أي اللابداية، والأبد، أي اللانهاية؛ وكلاهما مجرد ألفاظ، مجموعات حروف ذات دلالة ذاتية. ونستعمل في الإنجليزية \textenglish{Eternal}، من \textenglish{Æviternus}، أي الممتد عصرًا أو حياة، بالقدر نفسه من التساهل، فنطبّقها على ثلاث أفكار متميزة: العادي المتكرر في الاستعمال الشعبي، وما لا يخضع للمدة، وما يدوم أبدًا فيشمل كل مدة. و«صانع العلم المحيط» هو قول الشاك الروماني القديم: «صنع الإنسان الآلهة»، \textenglish{Homo fecit Deos}.

والقسم التالي نواح طويل واحد على تناقضات الوجود كله وأسراره، ونهايته المظلمة، وحزنه الذي لا يتناهى، وعلى سر الأسى الذي قال لوثر إنه يقوم تحت كل حياة. وكما عند يوربيديس: «أن تحيا هو أن تموت، وأن تموت هو أن تحيا». ويستعير الحاج عبدو الفكرة الهندوسية عن الجسد الإنساني. ويقول مانو: «إنه دار، عظامها عوارضها وروافدها، وأعصابها وأوتارها حبالها، وعضلاتها ودمها ملاطها، وجلدها غلافها الخارجي؛ لا تملؤها رائحة طيبة، بل تثقلها النجاسات؛ دار تجتاحها الشيخوخة والحزن، ومقر للمرض، تعذبه الأوجاع، وتطارده صفة الظلام، تاما گونا، ولا يستطيع الثبات». وبدأ الوعاء والفخّار عند المصريين القدماء: «يجلس كنيف في فيلة كفخّار عند دولابه، ويشكّل الطين، ويعطي روح الحياة لمنخري أوزيريس». ومن هنا «النفخة» في سفر التكوين. ثم نلقاه في الفيدات: الكائن «الذي يشكّل المزهرية الفخارية، والطين الذي تُصنع منه». ونجده بعد ذلك في قول إرميا: «قم وانزل إلى بيت الفخاري»، الفصل ١٨، الآية ٢؛ وأخيرًا في الرسالة إلى أهل رومية، الفصل ٩، الآية ٢٠: «أليس للفخاري سلطان على الطين؟» فلا عجب أن تكون اليد الأولى التي صاغت طين الإنسان فكرة مألوفة في الفكر الشرقي. و«قفر العذاب» بوذية أو شوبنهاورية خالصة بسيطة. وقد صغت «التراب على التراب» على قوافي «القديس إيزيدور» المعروفة، سنة ١٤٤٠:

\begin{SourceQuotation}
التراب من التراب صيغ صياغة عجيبة؛\\
والتراب من التراب نال كرامة من لا شيء؛\\
والتراب على التراب جعل كل فكره؛\\
كيف يُجلب إليه التراب على التراب، إلى آخره.
\end{SourceQuotation}

ويوحي «راكب الجمل» بقول أوسيان: «بضع سنوات أخرى، ثم تأتي عاصفة الصحراء». واختار العرب الجمل ذا السنام الواحد مركبًا للموت، ربما لأنه يحمل جثة البدوي إلى المقبرة البعيدة، حيث يرقد بين قرابته وأهله. وتذكّرنا نهاية هذا القسم بقول:

\begin{SourceQuotation}
ما أفقر الإنسان وما أغناه، ما أحطّه وما أرفعَه؛\\
ما أعقده وما أعجبَه!
\end{SourceQuotation}

وينتقل الحاج الآن إلى نتائج أفكاره الطويلة القلقة؛ وقد تعمدت تحوير استهلاله إلى صدى لميلتون:

\begin{SourceQuotation}
إلى أن تبلغ الخبرة العتيقة؛\\
شيئًا من نبرة النبوءة.
\end{SourceQuotation}

ويعلن بجرأة أن لا إله على النحو الذي خلق به الإنسان خالقه. وهنا يتفق مع الفكر الحديث. يقول ج ج روسو في الاعترافات، الكتاب الأول، الفصل السادس: «في الغالب يصنع المؤمنون الإله كما هم أنفسهم؛ فالطيّبون يجعلونه طيبًا، والأشرار يجعلونه شريرًا؛ والمتدينون الحاقدون المتعكرون لا يرون إلا الجحيم، لأنهم يريدون إدانة الجميع؛ أما النفوس المحبة اللطيفة فلا تكاد تؤمن به. ومن دهشاتي التي لا أفيق منها أن أرى فينيلون الطيب يتحدث عنه في تليماك كما لو أنه يؤمن به حقًا؛ لكنني آمل أنه كان يكذب حينها؛ فمهما كان المرء صادقًا، لا بد له أن يكذب أحيانًا حين يكون أسقفًا». ويقول شيلر: «يرسم الإنسان نفسه في آلهته». ولهذا أجاز إله الطبيعة، \textenglish{Naturgott}، الذي كان لجميع الشعوب القديمة وبدأ به كل مذهب، أفعالًا غير أخلاقية بوضوح، بل شيطانية كثيرًا ما تكون، واستحسنها. ولم يُضفَ على الاعتقاد طابع أخلاقي إلا حين بدأ ضمير الجماعة، ومعه ضمير أفرادها، يطمح إلى عصره الذهبي: الكمال. ويقول كارل فوغت: «الإله صيغة التفضيل التي صيغتها الأصلية الإنسان»؛ أي إن الفكرة الشعبية عن الألوهية هي إنسان مكبّر غير طبيعي.

ثم يقتبس مراجع حجته. فبوذا، الذي حولته الكنيسة الكاثوليكية إلى القديس يوسافات، رفض الاعتراف بإيشوارا، أي الألوهية، بسبب سر «قسوة الأشياء». وقد وجد شوبنهاور، متشائم الآنسة كوب النموذجي، الذي يمثل بوذا في عالم الفكر الغربي على أبعد مسافة متواضعة، رؤية تعاسة الإنسان، بصرف النظر عن أفعاله، طاغية إلى حد استنتج معه أن الإرادة العليا شريرة، «بلا قلب، جبانة، متغطرسة». وأنكر كونفوشيوس، «الملك بلا عرش، الأقوى من جميع الملوك»، إلهًا شخصيًا. وتسود الفكرة الأبيقورية صين الوقت الحاضر. ويقول السانتال الطورانيون في الهند الآرية: «الله عظيم، لكنه يقيم بعيدًا أكثر مما ينبغي»؛ وهذا هو اللسان العام للإنسان في الشرق الطوراني.

من الواضح أن الحاج عبدو يرى أن عبادة الأصنام تبدأ بإله شخصي. ولنلاحظ أن «المقالات التسع والثلاثين» تنفي هذا الأخير عمدًا. فالله عندها «كائن بلا أجزاء، أي بلا شخصية، ولا أهواء». وهو يعلن لاأدرية غامضة، ويعزو الإيمان الشعبي إلى أن «الخوف صنع الآلهة»، \textenglish{Timor fecit Deos}؛ فـ«كل دين، بلا استثناء، ابن الخوف والجهل»، كما يقول كارل فوغت. ويتكلم الآن بصفته «صابّ الخمر»، و«صاحب الحانة العتيق»، و«المجوسي العجوز»، و«شيخ المغان، أو المجوس»؛ وكلها ألقاب تُطلق على الصوفي في مقابل الزاهد. و«أصنامه» هي أوهام بيكون، التي «تقوم أسسها في تركيب الإنسان نفسه»، ولذلك تُدعى على نحو ملائم خرافات، \textenglish{fabulæ}. ومن السهل إثبات أن «سير الطبيعة المألوف» يقبل تفسيرات مختلفة. فقد كان أرسطو هادمًا عظيمًا كالإسكندر؛ لكن ابن اسطاغيرا شبه النبوي في العصور المظلمة، الذي حكم العالم حتى نهاية القرن الثالث عشر، صار عند مارتن لوثر «ملعونًا مرتين»؛ ثم ألغاه أخيرًا غاليليو ونيوتن. وقد حذفت هنا مقطعين. أولهما:

\begin{SourceQuotation}
نظريات بدل الحقائق، وخرافة بدل الواقع، ومذهب بدل العلم، تقلق الفكر؛\\
وتحتقرون درس الحياة العظيم الوحيد: أن تعرفوا أن كل ما نعرفه لا شيء.
\end{SourceQuotation}

وهذا في الواقع:

\begin{SourceQuotation}
أحسنت القول، يا أنبل أبناء أثينا؛\\
أقصى ما نعرفه أن لا شيء يمكن أن يُعرف.
\end{SourceQuotation}

والثاني:

\begin{SourceQuotation}
ماهية وجوهر، وتعاقب وعلة، وبداية ونهاية، ومكان وزمان؛\\
هذه لعب عقل الرجولة، مضحكة وسامية في آن واحد.
\end{SourceQuotation}

وليس هو الوحيد الذي ينظر هذه النظرة إلى «الزمان والمكان المزعجين». ومن تعريفات «اللامتناهي في الكبر» الحديثة أن فكرته تنشأ من نفي الشكل عن أي صورة؛ وفكرة «اللامتناهي في الصغر» من رفض المقدار لأي صورة. وهذا نموذج جيد لذلك «العلم الكئيب»: الميتافيزيقا.

ويقول مقطع آخر حُذف:

\begin{SourceQuotation}
كيف تدّعي، يا ظاهرة، أن تقيس الشيء في ذاته وتقدّر مداه؛\\
قل، أيهما أسهل فحصًا وإثباتًا: الوجود المطلق أم الإنسان الفاني؟
\end{SourceQuotation}

قد يظن المرء أنه قرأ كانط في «ما يمكن معرفته وما لا يمكن معرفته»، أو سمع بالسيدة الأمريكية التي استطاعت «التمييز بين المتناهي واللامتناهي». ومن أفكار العصر المألوفة في الغرب والشرق أن العلم محصور في الظواهر، ولا يستطيع بلوغ الأشياء في ذاتها، \textenglish{Noumena}. وهذه هي الواقعية المدرسية، «بقية ميتافيزيقا رديئة»، التي تشوّه مذهب كونت. ومع جميع دعاواها، فهي لا تعني إلا أن أشياء في ذاتها موجودة، أو يمكن تصورها، أي غير مرتبطة بالفكر؛ وأننا نعرف أنها موجودة، وفي الوقت نفسه لا نستطيع أن نعرف ما هي. لكن من يجرؤ أن يقول «لا يستطيع»؟ ومن يستطيع قياس عمل الإنسان حين يرتفع فوق ذواتنا الحاضرة بقدر ارتفاعنا فوق إنسان الكهف في الماضي؟

تشير «سلسلة الكون» إلى الفكرة الجاينية القائلة إن الكل، الذي يتكون من مبادئ عقلية وطبيعية، موجود منذ الأزل كله؛ وإنه خضع لدورات لا تنتهي، أسبابها القوى الأصيلة للطبيعة، العقلية والفيزيائية، بغير تدخل ألوهية. لكن الشاعر يسخر من «غير البشري»، أي غير أنفسنا، أي نفي أنفسنا، ومن ثم عدم الوجود. ويخلط معظم الشرقيين بين النقيضين، حيث يدل أحد الحدين على شيء والآخر على لا شيء، مثل أنفسنا وغير أنفسنا، وبين الضدين، مثل الغني وغير الغني، أي الفقير، حيث يعبّر كلا الحدين عن شيء. ولذلك فإن «اللامتناهي» الإيجابي السلبي ليس متمّم المتناهي، بل نفيه. ويسخر الغربي من هذه العملية بجعل «غير حصان» الكيان المكمّل لـ«حصان». ويختم الحاج بالمبدأ الصوفي المحبب أن الحواس الخمس، أم الست؟، هي أبواب كل معرفة إنسانية؛ وأنه ما من صورة للإنسان، أكانت تجسدًا للألوهية أم نبيًا أم رسولًا أم حكيمًا، أنتجت قط فكرة لم تتكون داخل دماغها بالعمل الوحيد لهذه العوامل المادية المبتذلة. ومن الواضح أنه ليس روحانيًا ولا مثاليًا.

ثم يبيّن أن الإنسان يرسم نفسه في إلهه، وأن «الإله تعبير عن العرق»: معلّم في النيل، وتجريد في الهند، ومنجّم في كلدان؛ حيث كان إبراهيم، كما يقول بيروسوس، نقلًا عن يوسيفوس، الآثار، الكتاب الأول، الفصل ٧، الفقرة ٢، والكتاب الثاني، الفصل ٩، الفقرة ٢، «ماهرًا في العلم السماوي». ويذكر أكرانا زمان، الزمن غير المنتهي، عند المجوس، والثنائي العامل، هرمزد وأهريمن. ويَسِم إله العبرانيين بالمحاربة والقسوة. وقد سمع بالمخلوقات الجميلة للخيال اليوناني الذي، إذ لم ينسب طبيعة أخلاقية إلى الألوهية، أدخل اللاهوت في علم الطبيعة؛ وبدا، مثل الأستاذ تندال، أنه يعدّ المادة كلها، في كل مكان، حية. وقد تبنّينا نحن واحدية مختلفة جدًا: اللاهوت بخالقه الواحد، ووحدة الوجود بـ«قوة الروح الواحدة المشكّلة»، والعلم بطاقته الواحدة. وهو شديد على المسيحية و«إلهها الثلاثي»؛ ولم ألطف عبارته، التي تعني «لغزًا»، وإن كان ذلك قد يؤذي القراء. فما من شيء أشد إلغازًا للعقل المسلم من التثليث المسيحي؛ يمكنه تجاوز جميع الاعتراضات الأخرى، لا هذا. وليس محبًا للإسلام أيضًا؛ فالإسلام، مثل المسيحية، له عبرانيته الزهدية وهيلينيته اللذّية، ويتحرك عالم الفكر بين هذين الطرفين. والأولى، التي تُعرّف بأنها الاهتمام الغالب أو الحصري بممارسة الصواب، يمثلها التأثير السامي والعربي، القرآني والحديثي. أما الثانية، دين الإنسانية، والشغف بالحياة والنور، والثقافة والذكاء، والفن والشعر والعلم، فيمثلها في الإسلام ذكرى ملوك المجوس وأبطالهم وجميلاتهم وشعرائهم وحكمائهم القدماء، تلك الذكرى المعزّزة بحب وقلة تقوى. ومن هنا ذكر زال وابنه رستم، وكورش، وكأس جمشيد، جامِ جمشيد، التي يمكن ترجمتها كأسًا أو مرآة؛ فقد كانت تُظهر العالم كله داخل حافتها، ولذلك سُمّيت جامِ جهان نما، كاشفة الكون. ومن الواضح أن العبارات الاحتقارية عن الغذاء من لبن الإبل ولحم السوسمار، أي السحلية الخضراء، مقتبسة من أبيات الفردوسي الشهيرة التي تبدأ بـ:

\begin{SourceQuotation}
بلغ بالعرب الأمر إلى حدّ...\\
\textenglish{Arab-râ be-jâî rasîd’est kâr}.
\end{SourceQuotation}

ويشتد الحاج على الذين يجعلون الألوهية «خوان يغما»، أو مائدة نهب، كما يقول الفرس. وهو يرى الرعاة رجالًا:

\begin{SourceQuotation}
يسلبون الخراف ليكسوا أنفسهم.
\end{SourceQuotation}

وهكذا يبيّن شوبنهاور بغضب، كما في كتاب «الحياة» وغيره لفيلهلم گيفينر، كيف «ينبغي للأمة الإنجليزية أن تعامل تلك الجماعة من المنافقين والدجالين وجامعي المال، أي رجال الدين الذين يلتهمون سنويًا ثلاثة ملايين وخمسمئة ألف جنيه».

ويقرر الحاج بعبارة صريحة أن لا خير ولا شر بالمعنى المطلق الذي صنعه الإنسان لهما. وهنا يتفق مع بوب:

\begin{SourceQuotation}
رغم الكبرياء، ورغم الطبيعة المخطئة؛\\
حقيقة واحدة واضحة: كل ما هو كائن صواب.
\end{SourceQuotation}

وللأسف فإن العكس صحيح بالقدر نفسه: كل ما هو كائن خطأ. والخضر هو إيليا الذي حيّر ميلمان. وهو يمثل الصوفي الباطني، بينما يمثل موسى الزاهد الظاهري؛ وقد استولى المسلمون على مغامراتهما الغريبة التي اخترعها اليهود. ويسخر من حرية إرادة الإنسان؛ ويكشف، مثل ديدرو، «مهرّجًا في أسقف، وساتيرًا في رئيس، وخنزيرًا في كاهن، ونعامة في وزير، وإوزة في كبير كتبة». ويتمسك بالحظ، تيخي عند ألكمان، أي المصادفة، أخت النظام والثقة، وابنة الاستبصار. وكانت غازلات القدر الإسكندنافيات أورد، أي ما كان، الماضي؛ وفيرداندي، أي ما يصير، الحاضر؛ وسكولد، أي ما سيكون، المستقبل. ويشير إلى أفلاطون الذي جعل الديميورغوس، الصانع، يخلق العوالم بواسطة اللوغوس، دابار بالعبرية، أي الكلمة الخالقة، عبر الأيونات. وكانت أيونات المتصوفة هذه فيوضًا روحية من أيون، أي، حرفيًا، موجة تدفق أو عصر أو مدة أو يوم؛ ومن هنا \textenglish{ævum} في اللاتينية، و\textenglish{Awen} في الويلزية، أي سيل الإلهام النازل على الشاعر. وجعل باسيليدس المصري المسيحي الخالق يُخرج سبعة أيونات أو امتلاءات؛ ومن اثنين منها، الحكمة والقوة، صدرت درجات الملائكة الثلاثمئة والخمس والستون. وكانت جميعها خاضعة لأمير للسماء يُدعى أبراكساس، وهو نفسه تحت توجيه الأيون الأكبر، الحكمة. ويمثل آخرون العلة الأولى على أنها أنتجت أيونًا، أو عقلًا خالصًا؛ والأول أنتج ثانيًا، وهكذا حتى العاشر. وكان هذا ماديًا بما يكفي للتأثير في الهيولى، التي اتخذت بذلك صورة روحية. وهكذا اجتمع المتنافيان في خطة الخلق.

وينكر عصور البوذيين الثلاثة: السعيد كله، والسعيد الممزوج بالشقاء، والشقي المشوب بالسعادة، وهو الحاضر. وكان للزرادشتيين أربعة عصور، كل منها ثلاثة آلاف سنة. وفي الأول حكم هرمزد، إله الخير، وحده؛ ثم بدأ أهريمن، إله الشر، يحكم تابعًا؛ وفي الثالث حكما متساويين؛ وفي الأخير الجاري الآن ظفر أهريمن.

وفي مواجهة الفكرة الشائعة بأن الإنسان سبّب شقاء هذا العالم، يستشهد بالعصور التي ولّد فيها الحجر الرملي الأحمر القديم أسماكًا هائلة تلتهم أمثالها؛ وأنتجت الصخور الأوليتية طغاة الزواحف العظيمة في الهواء والأرض والبحر؛ وهزّت فيها وحوش عصري الإيوسين والميوسين الأرض بخطاها الثقيلة. ولا يزال عالم المياه مشهدًا بشعًا للقسوة والمذبحة والدمار.

ويقرر أن الضمير مصادفة جغرافية وزمنية. وهكذا يجيب الفيلسوف الحديث الذي غمر روحه العجب والرهبة من شيئين: «السماء المرصعة بالنجوم فوقي، والقانون الأخلاقي في داخلي». ويجعل هذا الحس الأخير تطورًا للغرائز القطيعية والاجتماعية؛ ولذلك لاحظ الرحالة أن الأخلاقي آخر خطوة في التقدم العقلي. والموريون عنده هم الدانكالي الهمج وغيرهم من القبائل الزنوجية، الذين يقدمون كأس لبن بإحدى اليدين ويطعنون بالأخرى. وهو يترجم حرفيًا الكلمة الهندية هاتي، أي فيل: الحيوان ذو الهات، أي اليد أو الخرطوم. وأخيرًا يشير إلى عصر البراكين النشطة، وهو الحاضر الذي ليس إلا مؤقتًا، وإلى تبدّل القطب، والمشهد الذي سيُرى من المشتري، أو كوكب جوبيتر.

ويسأل الحاج مرة أخرى السؤال القديم، القديم: ما الحقيقة؟ ويجيب نفسه على طريقة إمبراطور الصين الحكيم: «ليس للحقيقة اسم لا يتغير». ويقول كاتب إنجليزي حديث: «لقد اقتنعت منذ زمن طويل، من خبرة حياتي بوصفي رائدًا لعقائد مخالفة كثيرة، تتحول سريعًا إلى عقائد مقبولة، بأن الحقيقة كلها تقريبًا مزاجية بالنسبة إلينا، أو معطاة في العواطف والحدوس؛ وبأن النقاش والبحث لا يفعلان أكثر من تغذية المزاج». ويبدو أن شاعرنا يعني أن قوى الإدراك، حين تدرك إدراكًا صادقًا، تنقل حقيقة موضوعية عامة؛ بينما لا يمنح التأمل والمشاعر، أي عمل المنطقة الأخلاقية أو الفص الأوسط عند أصحاب علم فراسة الجمجمة، إلا حقيقة ذاتية شخصية فردية. وهكذا تمثل الحكمة «الثروات محرضات الشرور»، \textenglish{Opes irritamenta malorum}، واقعة متميزة عند إنسان؛ بينما يرى آخر أن الثروة حافز إلى الخير. ومن الواضح أن كليهما على صواب بحسب ما لديه من نور.

ويستشهد الحاج عبدو بأفلاطون وأرسطو، كما يفعل الشعراء الشرقيون عادة، وهم يستمتعون بالمنطق. ويبدو أنه يعني هنا أن القضية الكاذبة قضية واقعية بقدر القضية الصادقة. ومن الواضح أن «الإيمان ينقل الجبال»، و«يبقى الإيمان ثابتًا»، \textenglish{Manet immota fides}، اقتباسان. ويسخر من تعليم مجمع الفاتيكان الأول، الفصل الخامس، بأن «المؤمنين جميعًا أطفال صغار يصغون إلى صوت القديس بطرس»، الذي هو «أمير الرسل». ويشير إلى خيال بعض الفيزيائيين المحدثين أن «التقوى تغير جزيئي محدد في تلافيف اللب الرمادي». ويذكر بازدراء اللغز الذي يتحدث عنه ميلتون بتلك السلاسة، حيث كان المتحاورون:

\begin{SourceQuotation}
يستدلون استدلالًا عاليًا في العناية الإلهية والعلم السابق، والإرادة والقدر؛\\
والقدر الثابت والإرادة الحرة والعلم السابق المطلق.
\end{SourceQuotation}

وفي مقابل العقائد المحمدية السائدة التي تجعل نفس الإنسان أناَه المدركة، كيانًا ووحدة، يعدّها الصوفي خيالًا في مقابل الجسد الذي هو واقع؛ أو، في أقصى الأحوال، حالًا من الأحوال لا شيئًا؛ واتفاقًا بين قوى ليست أجسامنا إلا ظواهرها. ولا يخالف هذا أسطورة سفر التكوين. فـ«رُواح» العبرية و«روح» العربية، اللتان حُرّف معناهما الآن إلى النفس أو الروح، تعنيان ببساطة الريح أو النفَس، العلامة الخارجية المرئية للحياة. ومدارسهم المتأخرة أوضح أيضًا: «ما يحدث للإنسان يحدث للبهائم؛ كما يموت هذا يموت ذاك؛ ولهم جميعًا موت واحد، ويذهب الجميع إلى مكان واحد»، سفر الجامعة، الفصل ٣، الآية ١٩. لكن النفس الحديثة، وهي لا شيء وسلسلة من النفي، ونفي في المقام الأول، تُوصف في المهابهاراتا هكذا: «غير قابلة للتقسيم، ولا للتصور، ولا للإدراك؛ أبدية وعامة ودائمة وثابتة؛ غير مرئية وغير قابلة للتبدل». ومن هنا الروحانية الحديثة التي، وهي ترفض المادية، لا تستطيع استعمال إلا لغة مادية.

ويقول الحاج إن هذه مجرد أصوات. وهو لا يقرر «الكلمات تولّد الكلمات»، \textenglish{Verba gignunt verba}، بل «الكلمات تولّد الأشياء»، \textenglish{Verba gignunt res}، خطوة أبعد. والفكرة هي «أصنام السوق، أشدها إزعاجًا» عند بيكون، \textenglish{idola fori, omnium molestissima}: أوهام اللغة ذات الوجهين؛ إما أسماء أشياء ليس لها وجود في الواقع، وإما أسماء أشياء فكرتها مضطربة سيئة التحديد.

ويشتق فكرة الروح من «الشبح الهمجي» الذي عرّفه الدكتور جونسون بأنه «نوع من كائن ظلي». وهو يلاحظ بحق أنها نشأت، ربما، في مصر، ولم يخترعها «أهل الكتاب». ويقصد المسلمون بهذا اللفظ اليهود والمسيحيين الذين لهم وحي معترف به، بينما يرفض جهلهم الاعتراف به للمجوس والهندوس والكونفوشيوسيين.

ومن الواضح أنه يتمسك بمذهب التقدم. فالبروتوبلازم عنده هو الإيلياسترون، والمادة الأولى، \textenglish{Prima Materies}. وكلمتنا \textenglish{matter} مشتقة من السنسكريتية ماترا، التي تعني، على الوجه الصحيح، المثال غير المرئي للمادة المرئية؛ أو، باللغة الحديثة، الجوهر المتميز عن مجموع خصائصه الفيزيائية والكيميائية. ومن ثم لا توجد ماترا إلا في الفكر، ولا يمكن التعرف إليها بفعل الحواس الخمس. وتذكّرنا «سلسلة الوجود» عنده بشجرة نسب الحصان للأستاذ هكسلي: أوروهيبوس، وميزوهيبوس، وميوهيبوس، وبروتوهيبوس، وبليوهيبوس، وإيكوس. ومن الواضح أنه سمع بعلم الأحياء الحديث، أو مذهب حياة المادة، الذي يعد أنواعه البالغة ربع مليون نوع من الكائنات الحية، الحيوانية والنباتية، تعديلات متدرجة لوحدة أساسية عظيمة؛ وحدة لما يُسمّى «القوى العقلية» كما للبنية الجسدية. وهذه هي بقعة الهلام. ويسخر من الفكرة الشائعة أن الإنسان هو الشخصية المركزية العظمى التي يدور حولها كل شيء كالدمى؛ أي حقبة درابر المتمركزة حول الإنسان، التي تعيش، ويا للعجب، إلى جانب التلسكوب والمجهر. وما دام الإنسان حديث الأصل، وقد ينتهي في حقبة مبكرة من العالم الأكبر، فإن جميع الأشياء قبل ميلاده كانت تدور حول لا شيء، وقد تستمر في ذلك بعد موته.

ويرى الحاج، الذي يندد في موضع آخر بـ«الجهل المركب»، أن كل شر يأتي من الخطأ؛ وأن كل معرفة تطورت بهدم الخطأ، وهو المجرى المعتاد للفكر الإنساني. ويختم هذا القسم بحقيقة عظيمة: ثمة أشياء لا يستطيع العقل الإنساني، أو الغريزة الناضجة، في حاله غير المتطورة، أن يحيط بها؛ لكن العقل ناموس لنفسه. ولذلك لسنا ملزمين بأن نؤمن، أو نحاول الإيمان، بأي شيء مخالف للعقل أو مناقض له. وهنا يعارض روما معارضة مباشرة؛ فهي تقول: «لا تحتج بالتاريخ؛ فهذا حكم خاص. ولا تحتج بالكتاب المقدس؛ فهذه هرطقة. ولا تحتج بالعقل؛ فهذه عقلانية».

ويرى، مع آباء الكتاب المقدس العبري، أن الحياة الحاضرة كافية كل الكفاية لكائن عقلي، لا عاطفي؛ ولذلك لا حاجة إلى جنة أو جحيم. أما الشاعر الغربي فيغني بتناقض أكبر بكثير:

\begin{SourceQuotation}
ليس للجحيم حدود، ولا ينحصر في موضع واحد بعينه؛\\
بل حين نكون في الجحيم، وحيث يكون الجحيم، يجب أن نظل دائمًا؛\\
وباختصار، حين ينحل هذا العالم كله ويتطهر كل مخلوق؛\\
ستكون جميع الأمكنة التي ليست جنة جحيمًا.
\end{SourceQuotation}

فأي حاجة إلى جحيم حين يكون الجميع أطهارًا؟ وهو يوسّع النظرية البوذية القديمة القائلة إن السعادة والشقاء موزعان بالتساوي بين الناس والبهائم؛ فبعضهم يتمتع كثيرًا ويتألم كثيرًا، وآخرون بالعكس. ومن هنا يقرر ديدرو: «الأهواء المعتدلة لا تنتج إلا المألوف... والرجل المعتدل الهوى يعيش ويموت كبهيمة». ومرة أخرى نجد نصف الحقيقة:

\begin{SourceQuotation}
علامة الرتبة في الطبيعة؛\\
هي القدرة على الألم.
\end{SourceQuotation}

وتتضمن هذه الأخيرة قدرة مساوية على اللذة؛ وهكذا يُحفظ التوازن.

ثم يبيّن الحاج عبدو أن الإيمان مصادفة الميلاد. ويقول أحد أبياته المزدوجة المحذوفة:

\begin{SourceQuotation}
العرق يصنع الدين؛ صحيح! لكن المصنوع يؤثر دائمًا في الصانع؛\\
وإله متناهٍ، وخطيئة لا متناهية، يحطّان الإنسان بدل رفعه.
\end{SourceQuotation}

ويُدخل في صورة حوار شعوبًا مختلفة، يقاتل كل منها لإثبات اعتقاده الخاص. فيسيء الفرنجي المسيحي إلى الهندي، ويرد هذا بأنه من دم الملتشا، أي المختلط أو غير النقي، وهي كلمة تُطلق على غير الهندوس جميعًا. ويفعل الشيء نفسه الناصري والمحمدي؛ والكونفوشيوسي الذي لا يؤمن بشيء؛ والصوفي الذي تكون له، بطبيعة الحال، الكلمة الأخيرة. وربط العذراء مريم والقديس يوسف بالثالوث في الكنيستين الرومانية واليونانية يجعل كثيرًا من المسلمين يستنتجون أن المسيحيين لا يؤمنون بثلاثة أقانيم، بل بخمسة. ويكتب إنجليزي عن الآباء الأوائل: «لم يقولوا فقط إن ثلاثة تساوي واحدًا وواحدًا يساوي ثلاثة؛ بل ادّعوا تفسير كيف نشأ هذا التركيب الحسابي العجيب. فقد قُسم ما لا يقبل القسمة، ومع ذلك لم يُقسم؛ وكان قابلًا للقسمة، ومع ذلك غير قابل لها؛ وكان الأسود أبيض والأبيض أسود، ومع ذلك لم يكن ثمة لونان بل لون واحد؛ وكل من لا يؤمن بذلك سيُدان». ويجري الاقتباس العربي في الأصل هكذا:

\begin{SourceQuotation}
أحسنُ المكانِ للفتى الجهنّمُ.\\
أفضل الأمكنة للفتى الكريم هو جهنم.
\end{SourceQuotation}

وجهنم، أو الجحيم، هي الموضع الناري للعقاب الأبدي. وكذلك يُدان القول الثاني «النار ولا العار» عند المتمسك بالقرآن. و«وقاحة» القدر، گوستاخي، هي عبارة عمر الخيام في المقطع الثلاثين:

\begin{SourceQuotation}
ماذا؟ من أين سِيق بنا إلى هنا بغير سؤال؛\\
وإلى أين يُساق بنا من هنا بغير سؤال؛\\
آه، كم من كأس من هذه الخمر المحرمة؛\\
يلزم لإغراق ذكرى تلك الوقاحة!
\end{SourceQuotation}

أما في المعنى الصوفي، فالكلمة تعني «دلال المحبوب»، تلك «الجزئية من النفخة الإلهية»، \textenglish{divinæ particula auræ}. وينتهي القسم بقول بوب:

\begin{SourceQuotation}
لا يمكن أن يخطئ من كانت حياته على صواب.
\end{SourceQuotation}
